\documentclass{cumcmthesis}
% 电子版提交时通常可使用：\documentclass[withoutpreface]{cumcmthesis}

\usepackage{booktabs}
\usepackage{amsmath,amssymb,bm}
\usepackage{subcaption}
\usepackage{graphicx}
\usepackage{float}
\usepackage{algorithm}
\usepackage{algorithmic}
\usepackage{cleveref}
\usepackage[most]{tcolorbox}
% cumcmthesis v2.6 自定义了旧版 \supercite；在加载 biblatex 前释放该命令。
\let\supercite\relax
\usepackage[backend=biber,style=numeric,sorting=none]{biblatex}
\addbibresource{references.bib}

% 命题按阿拉伯章节号编号。0--4 是默认正文阅读预算，不是数学硬上限；超预算时应有实际建模作用与必要性说明。
% 短证明与命题正文放在同一不可分页外框；技术细节过长时移到附录。
\newtheorem{proposition}{命题}[section]
\renewcommand{\theproposition}{\arabic{section}.\arabic{proposition}}
\newenvironment{hskproposition}[1]
  {\begin{tcolorbox}[
      colback=white,
      colframe=black!72,
      boxrule=0.65pt,
      arc=1mm,
      outer arc=1mm,
      left=8pt,
      right=8pt,
      top=7pt,
      bottom=7pt,
      before skip=8pt,
      after skip=8pt,
      boxsep=0pt
    ]
    \begin{proposition}[#1]}
  {\end{proposition}\end{tcolorbox}}
\newenvironment{hskproof}
  {\par\smallskip\noindent\textbf{证明：}\par\nopagebreak\smallskip}
  {\hfill$\square$\par}

% 按当届题面和队伍信息修改，不在模板中固定年份、题号和小问数。
\newcommand{\HSKProblemCode}{A}
\newcommand{\HSKTeamNumber}{0000}
\newcommand{\HSKContestYear}{YYYY}
\newcommand{\HSKContestMonth}{9}
\newcommand{\HSKContestDay}{1}

\title{XXX问题的XXX建模与分析}
\tihao{\HSKProblemCode}
\baominghao{\HSKTeamNumber}
\schoolname{学校名称}
\membera{成员A}
\memberb{成员B}
\memberc{成员C}
\supervisor{}
\yearinput{\HSKContestYear}
\monthinput{\HSKContestMonth}
\dayinput{\HSKContestDay}

\begin{document}
\maketitle

\begin{abstract}
% 第一段通常 2--3 句：研究对象 + 核心矛盾 + 全文统一路线。不要在第一段堆全部小问方法清单。
……

% 随后按问写“模型/机制 + 决定性数值 + 直接结论”，避免每段都复制“针对问题X，本文……”。
% 优先正向陈述“建立了什么—得到什么—说明什么”，真实边界按需补充。
……

\keywords{关键词1\quad 关键词2\quad 关键词3\quad 关键词4}
\end{abstract}

\section{问题重述}
\subsection{问题背景}
% 默认一个自然段。优先利用题目给出的现实背景和领域矛盾，自然收束到待解决的问题。
% 不另起一段介绍“全文结构”“后文安排”，也不机械堆宏大时事背景。

\subsection{问题提出}
% 用自己的理解逐问转述“研究对象 + 关键条件 + 待求输出”，不提前写模型名、公式、算法和结果。
\noindent\textbf{问题一：}……

\noindent\textbf{问题二：}……

% 按题目实际小问继续，不固定五问。

\section{问题分析}
\subsection{问题一的分析}
% 通常一个自然段：难点、对象关系、跨问依赖、建模抓手。正式公式和最终结果默认留到模型/结果章节。

\subsection{问题二的分析}
% 按实际小问继续；不要写“先清洗、再建模、再求解、再画图”的流水线。

\section{模型假设}
% 正文只使用自然序号，不显示 H1/H2、A1/A2 等内部合同编号。
\begin{enumerate}
    \item ……。
    \item ……。
    \item ……。
\end{enumerate}

\section{符号说明}
% 只列全文高频符号。场景/模型/方案等优先考虑短上标，避免长文本和多层下标。
\begin{table}[H]
    \centering
    \caption{主要符号说明}
    \label{tab:symbols}
    \begin{tabular}{cccc}
        \toprule
        符号 & 含义 & 单位 & 类型 \\
        \midrule
        $x_i$ & 元素/位置索引示例 & -- & 变量 \\
        $I^{(s)}$ & 第 $s$ 个模型对应结果 & -- & 状态量 \\
        $\theta$ & 参数示例 & -- & 参数 \\
        \bottomrule
    \end{tabular}
\end{table}

% preprocessing_decision=not_needed 时删除本章，只在数据说明中简洁交代数据审计结论。
% question_local 时把必要变换放入对应问题；project_level 时才保留实质数据预处理章节。
\section{数据预处理}
\subsection{数据问题与处理必要性}
% 说明真实数据问题；不得为了形式完整编造清洗、插值、标准化。

\subsection{处理方法与验证}
% project_level 时给公式、参数依据、验证、处理前后证据和后续模型接口。

\section{问题一模型建立及求解}
\subsection{模型建立}
% 从本题对象和机制开始，详细推导关键关系。
设……，得到
\begin{equation}
    F(x,\theta)=0,
    \label{eq:q1-core-relation}
\end{equation}
式~\eqref{eq:q1-core-relation} 控制……，并用于下一步构造……。

% 命题只在确实承担降维、等价变换、单调性、可行性等建模作用时使用。
% 短证明默认自然分段；只有明显多阶段证明才使用 enumerate。
% \begin{hskproposition}{变换模型的等价性}
% \label{prop:q1-equivalence}
% 在条件……及参数范围……内，原模型与变换模型具有相同的可行解映射和最优目标值。
% \begin{hskproof}
% 由定义构造……，在当前定义域内该映射一一对应。又因为……，约束在映射前后保持一致。
%
% 进一步比较目标函数可得……，故两个模型的最优值与最优解映射一致，命题得证。
% \end{hskproof}
% \end{hskproposition}
% 框后另起段说明命题对当前模型的实际作用和失效边界。

% 核心模型收束按项目框架中的 required / inline / not_applicable 自适应。
% 当前示例展示 required 形态；若为 inline，则把最终一两个计算/判定关系留在模型建立末尾并删除本小节；
% 若为 not_applicable（直接计算、结果读取或完全继承且无新增模型），删除本小节。
\subsection{核心模型汇总}
\begin{equation}
\left\{
\begin{aligned}
\min\quad & f(x) \\
\text{s.t.}\quad & g_i(x)\le 0, \quad i=1,\ldots,m, \\
& h_j(x)=0, \quad j=1,\ldots,n, \\
& x\in\Omega .
\end{aligned}
\right.
\label{eq:q1-model-summary}
\end{equation}
% 非优化题按本题结构改成“状态方程—观测关系—概率映射—边界/初值—输出量”的收束形式。

\subsection{模型求解}
% 写实际算法、初始化/参数、停止条件、可行性/收敛口径；不要只写“利用Python/MATLAB求解”。

\subsection{求解结果}
% 主结果、方案和核心数值放在这里。正文核心图/表必须有邻近的显式交叉引用，例如：
% 图~\ref{fig:q1-main} 给出了……；从图中的……可以看到……，这与式~\eqref{eq:q1-core-relation}中的……一致。
% 表~\ref{tab:q1-main} 汇总了……；关键差异主要集中在……。
% 引用方式要自然变化，不要每段都写“由图可知”。

% 若存在实质深化证据，按题目专属名称继续设置小节，例如：
% \subsection{阈值与失效边界}
% \subsection{结构稳健性}
% 最后一个结果/深化段自然回答本问，不默认再建“问题一结论”。

% 按实际题目继续设置问题二及后续；每问的核心模型收束状态都由当前模型复杂度决定，不机械复制同名小节。
% 共享基础模型只完整定义一次，后问只写新增变量、目标、约束和证据。

\section{模型的评价与推广}
% 若确有实质“改进”内容，可改名为“模型的改进、评价与推广”。
% “模型的改进”和“模型的推广”均为选写，只保留确实有内容的部分。
\subsection{模型的优点}
% 数量由实际模型决定；每条必须对应模型结构、结果或验证证据，不设置优缺点条目的机械数量关系。
\begin{enumerate}
    \item ……。
    \item ……。
    \item ……。
\end{enumerate}

\subsection{模型的缺点}
% 数量由实际限制决定；写清限制会影响哪个结果、参数范围或适用边界。
\begin{enumerate}
    \item ……。
    \item ……。
\end{enumerate}

\subsection{模型的推广}
% 选写。简要说明可迁移对象、需要替换的参数/边界和不能直接迁移的条件。

% 中文国赛默认不设置独立“结论”一级章。各问直接答案留在对应“求解结果”及深化分析中。
% 只有当届模板、用户或论文类型明确要求时才新增结论章。

\printbibliography[title={参考文献}]

\begin{appendices}
\section{核心代码与复现说明}
% 正文只保留必要伪代码；完整代码可放附录或独立附件，并说明输入、功能和输出。

\section{补充推导与图表}
% 技术证明超过正文安全长度时放在此处；正文命题框只保留关键证明链。
% 仅放正文不宜展开但仍有证据价值的图表。
\end{appendices}

\end{document}